% Chapter 2: the news summary report split into five modules, with the data
% (interfaces) passed between them.
% Shared preamble for the book's diagrams.
%
% Each figures/<name>.tex is a standalone TikZ picture. `make figures` builds,
% for each language <lang> (ko, en),
%   img/<lang>/<name>.pdf  vector, used by the print book (book.lua swaps it in)
%   img/<lang>/<name>.png  300dpi, white background, used by the Markdown
%
% Text is written once per language as \tr{한국어}{English}. The language is
% chosen by defining \figlang before this file is read, e.g.
%   xelatex -jobname=<name> "\def\figlang{en}\input{figures/<name>}"
% and defaults to ko.
%
% Figures are drawn to fit the print text block (95mm wide) at 1:1 scale, in
% grayscale only so they survive black-and-white printing.
\documentclass{article}
\providecommand\figlang{ko}
\def\figlangKo{ko}
\ifx\figlang\figlangKo
  \newcommand\tr[2]{#1}
\else
  \newcommand\tr[2]{#2}
\fi
\usepackage{fontspec}
\ifx\figlang\figlangKo
  \usepackage[hangul]{xetexko}
  \setmainfont{Noto Sans CJK KR}
  \setmainhangulfont{Noto Sans CJK KR}
\else
  \setmainfont{Noto Sans}
  % No hyphenation: diagram labels read better with whole words.
  \AtBeginDocument{\hyphenpenalty=10000 \exhyphenpenalty=10000 }
\fi
\usepackage{xcolor}
\usepackage{tikz}
\usetikzlibrary{arrows.meta,positioning,calc,fit,decorations.pathmorphing,decorations.pathreplacing,shapes.geometric,backgrounds}

% Grays (print-safe)
\definecolor{ink}{gray}{0.12}   % lines and text
\definecolor{mid}{gray}{0.45}   % secondary text, guides
\definecolor{soft}{gray}{0.72}  % light strokes
\definecolor{lite}{gray}{0.90}  % filled shapes (sticky notes, boxes)
\definecolor{pale}{gray}{0.96}  % large background areas

% Font sizes in pt, relative to the 9.5pt body text.
\newcommand\figtext{\fontsize{7.5}{10}\selectfont}
\newcommand\figsmall{\fontsize{6.5}{8.5}\selectfont}
\newcommand\figtiny{\fontsize{5.5}{7}\selectfont}

\tikzset{
  every picture/.style={line width=0.6pt, draw=ink, text=ink, font=\figtext, >={Stealth[length=2.2mm,width=1.6mm]}},
  arr/.style={->, draw=ink, line width=0.6pt},
}
\setlength\parindent{0pt}
\pagestyle{empty}
% Break Korean lines only at spaces (between words), not between syllables,
% which reads better in narrow diagram boxes.
\ifx\figlang\figlangKo
  \AtBeginDocument{\XeTeXlinebreakpenalty=10000 }
\fi

% Crop each page to its picture (+1mm border); TeX Live here has no
% standalone/preview packages. Wrap the tikzpicture in \begin{fig}...\end{fig}.
\newsavebox\figbox
\newenvironment{fig}{\begin{lrbox}{\figbox}}{\end{lrbox}\shipfig}
\newcommand\shipfig{%
  \hoffset=-1in \voffset=-1in
  \pdfpagewidth=\dimexpr\wd\figbox+2mm\relax
  \pdfpageheight=\dimexpr\ht\figbox+\dp\figbox+2mm\relax
  \shipout\vbox{\kern1mm\hbox{\kern1mm\usebox\figbox}\kern\dp\figbox\kern1mm}}

\begin{document}\begin{fig}
\begin{tikzpicture}[
    mod/.style={draw=ink, line width=0.8pt, fill=lite, rounded corners=1mm,
                minimum width=30mm, minimum height=7mm, font=\figtext\bfseries, align=center},
    dat/.style={draw=mid, rounded corners=3mm, minimum width=30mm, minimum height=6mm,
                font=\figsmall, align=center},
    note/.style={font=\figsmall, text=mid, align=left, text width=50mm, anchor=west},
    node distance=3.2mm,
  ]
  \node[dat] (in) {\tr{설정 키워드 · 오늘 날짜}{keywords · today's date}};
  \node[mod, below=of in] (search) {\tr{검색}{Search}};
  \node[dat, below=of search] (links) {\tr{뉴스 링크들}{news links}};
  \node[mod, below=of links] (fetch) {\tr{링크 콘텐츠 수집}{Fetch link content}};
  \node[dat, below=of fetch] (contents) {\tr{링크 콘텐츠들}{link contents}};
  \node[mod, below=of contents] (sum) {\tr{콘텐츠 요약}{Summarize}};
  \node[dat, below=of sum] (sums) {\tr{요약된 콘텐츠들}{summaries}};
  \node[mod, below=of sums] (fmt) {\tr{리포트 포맷터}{Report formatter}};
  \node[dat, below=of fmt] (rep) {\tr{리포트}{report}};
  \node[mod, below=of rep] (send) {\tr{리포트 전달}{Deliver report}};
  \node[dat, below=of send, font=\figsmall\bfseries] (inbox) {\tr{사용자의 메일함}{user's inbox}};

  \foreach \a/\b in {in/search, search/links, links/fetch, fetch/contents,
                     contents/sum, sum/sums, sums/fmt, fmt/rep,
                     rep/send, send/inbox}
    \draw[arr] (\a) -- (\b);

  \node[note] at ([xshift=5mm]search.east) {\tr{키워드와 날짜로 관련 뉴스 링크를 찾는다}{finds news links for the keywords and date}};
  \node[note] at ([xshift=5mm]fetch.east) {\tr{링크마다 본문을 가져온다}{gets the body of each link}};
  \node[note] at ([xshift=5mm]sum.east) {\tr{본문마다 짧게 요약한다}{writes a short summary of each body}};
  \node[note] at ([xshift=5mm]fmt.east) {\tr{요약들을 정해진 형식으로 묶는다}{puts the summaries into a set format}};
  \node[note] at ([xshift=5mm]send.east) {\tr{리포트를 사용자에게 메일로 보낸다}{emails the report to the user}};
\end{tikzpicture}
\end{fig}\end{document}
